\chapter{Подготовка данных и классификация банковских операций}

\section{Постановка задачи и обоснование подхода}

Банковская выписка является наиболее полным и достоверным первичным источником информации о денежных потоках предприятия: в отличие от управленческого учёта, она не содержит начислений без оплаты и отражает реальное движение денег на счёте. Вместе с тем сырая выписка бесполезна для финансового анализа — строки вида «оплата по счёту 1254 от ООО Прогресс» не несут управленческой информации о природе платежа.

Задача классификации банковских транзакций по экономическому содержанию является частным случаем задачи категоризации коротких текстов \cite{aggarwal2012survey}. В классической постановке она решается методами обработки естественного языка (NLP): наивный байесовский классификатор, метод опорных векторов на TF-IDF признаках \cite{sebastiani2002machine}, а в последние годы — предобученные трансформеры типа BERT \cite{devlin2018bert}.

Однако в контексте малого бизнеса задача имеет несколько особенностей, которые определяют выбор методов:
\begin{enumerate}
  \item \textbf{Высокая доля структурированных паттернов.} Назначения налоговых платежей, банковских комиссий и кредитных платежей имеют почти стандартизованные формулировки, поддающиеся регулярным выражениям.
  \item \textbf{Небольшой объём уникальных контрагентов.} В выборке из 4 086 операций число уникальных контрагентов ограничено; многие из них поддаются однозначной идентификации по имени.
  \item \textbf{Отсутствие размеченного обучающего набора.} Для обучения полноценного ML-классификатора необходима разметка, которой изначально нет.
  \item \textbf{Требование интерпретируемости.} Финансовый аудит требует проверяемости каждого решения — автоматическая классификация должна объяснять своё решение.
\end{enumerate}

Исходя из этого, в проекте принята многоуровневая стратегия: \textit{rule-based} классификация как основной метод, \textit{LLM-разметка} как вспомогательный сигнал для неоднозначных случаев, и ручная валидация выборки для оценки качества. Такой подход соответствует рекомендациям по применению гибридных методов классификации в задачах финансового мониторинга \cite{liu2004financial}.

\section{Источники данных и структура банковской выписки}

Исходные данные получены из двух источников: банковские выписки по расчётному счёту предприятия (формат Excel, экспорт из интернет-банка) и данные о продажах и оплатах из системы 1С (форматы Excel и CSV).

Банковская выписка за период май 2025 — апрель 2026 содержала несколько файлов с пересекающимися периодами. До очистки суммарный объём строк превышал 4 500. После обработки получена единая таблица из \textbf{4 086 транзакций}.

Структура итоговой таблицы банковских операций (\texttt{bank\_statement\_transactions}) включает:
\begin{itemize}
  \item \texttt{operation\_date} — дата операции;
  \item \texttt{document\_number} — номер платёжного документа;
  \item \texttt{direction} — тип операции: \texttt{credit} (поступление) или \texttt{debit} (списание);
  \item \texttt{amount} — сумма операции;
  \item \texttt{counterparty\_name} — наименование контрагента;
  \item \texttt{counterparty\_inn} — ИНН контрагента;
  \item \texttt{payment\_purpose} — назначение платежа (свободный текст);
  \item \texttt{source\_*} — служебные поля происхождения строки для аудита.
\end{itemize}

Для каждого файла при загрузке рассчитывается SHA-256 хэш, что позволяет реализовать инкрементальную логику пайплайна: уже обработанные файлы пропускаются, при изменении файла старые строки удаляются и файл парсится заново. Дополнительно пайплайн сверяет суммарные обороты по credit и debit со значениями из шапки банковской выписки; несовпадение приводит к остановке с ошибкой — это гарантирует целостность данных.

\section{Очистка данных}

\subsection{Дедупликация и пересекающиеся периоды}

Банк формирует выписку помесячно, и при склейке нескольких файлов операции на границах периодов могут задублироваться. Стандартная дедупликация по комбинации (дата, контрагент, сумма) неприменима: два платежа одному поставщику на одинаковую сумму в один день могут быть разными реальными операциями.

В проекте применена стратегия дедупликации по полному набору полей: \texttt{operation\_date}, \texttt{document\_number}, \texttt{counterparty\_inn}, \texttt{amount}, \texttt{payment\_purpose}. Технические дубли — строки с идентичными значениями всех перечисленных полей, попавшие в базу из пересекающихся файлов, — удаляются. Это позволяет сохранить реальные повторяющиеся операции (например, еженедельные арендные платежи одного размера).

\subsection{Парсинг дат}

Банковские выгрузки в формате Excel нередко содержат даты в смешанных форматах или с артефактами: дата может быть представлена как числовой код («серийный номер» Excel), строка в нестандартном формате или \texttt{NaT} из-за некорректного чтения файла.

Специализированная функция \texttt{parse\_bank\_dates} последовательно применяет несколько стратегий распознавания: преобразование числового кода Excel через \texttt{xlrd.xldate\_as\_datetime}, разбор строки по шаблонам \texttt{DD.MM.YYYY} и \texttt{YYYY-MM-DD}, fallback к контекстной интерполяции. После применения функции число операций без корректно распознанной даты снизилось до \textbf{нуля}.

\subsection{Нормализация назначений платежей}

Назначение платежа — свободный текст, заполняемый отправителем. Разные контрагенты описывают одинаковые по экономическому содержанию операции по-разному: «оплата по дог. б/о», «платёж согл. договору без объекта», «по счёту 1254», «за сырьё согл. дог. №12 от 01.01.25». Для rule-based классификации необходима нормализация:
\begin{itemize}
  \item приведение к нижнему регистру;
  \item удаление лишних пробелов и специальных символов;
  \item унификация частых аббревиатур («дог.» → «договор», «б/о» → «без объекта»);
  \item выделение токенов по ключевым паттернам.
\end{itemize}

\section{Разработка справочника категорий}

Разработка строгого справочника категорий — один из ключевых этапов проекта. Без чёткой таксономии автоматическая классификация лишена смысла: неоднозначные границы между категориями приводят к систематическим ошибкам, которые нельзя исправить никакими улучшениями алгоритма.

В проекте разработан справочник из 17 укрупнённых категорий без смысловых пересечений. Принципы построения:

\begin{enumerate}
  \item \textbf{Взаимная исключительность.} Каждая операция может относиться ровно к одной категории.
  \item \textbf{Соответствие типу операции.} Для каждой категории жёстко задан допустимый тип: \texttt{credit} (поступление) или \texttt{debit} (списание). Например, «оплаты от клиентов» — только \texttt{credit}, «производственные закупки» — только \texttt{debit}.
  \item \textbf{Разделение категорий и подкатегорий.} Детализация (конкретный вид сырья, тип налога, наименование банка-кредитора) выносится в поле \texttt{subcategory}, не смешивается с уровнем категории.
  \item \textbf{Управленческая интерпретируемость.} Итоговые категории должны напрямую соответствовать строкам управленческого отчёта о движении денежных средств.
\end{enumerate}

Полный справочник категорий приведён в таблице \ref{tab:categories}.

\begin{table}[H]
\centering
\caption{Справочник категорий банковских операций}
\label{tab:categories}
\begin{tabular}{llc}
\toprule
\textbf{Категория} & \textbf{Cash flow секция} & \textbf{Тип} \\
\midrule
Оплаты от клиентов & operating & credit \\
Производственные закупки & operating & debit \\
Упаковка & operating & debit \\
Аренда и коммунальные платежи & operating & debit \\
Персонал & operating & debit \\
Подрядчики и профессиональные услуги & operating & debit \\
Налоги и страховые взносы & operating & debit \\
Банковские комиссии & operating & debit \\
Цифровые сервисы, ПО и связь & operating & debit \\
Логистика & operating & debit \\
Маркетинг & operating & debit \\
Прочие операционные расходы & operating & debit \\
Возвраты от покупателей & returns & credit \\
Возвраты поставщикам & returns & credit \\
Кредиты и займы & financing & credit/debit \\
Внутренние переводы & internal & credit/debit \\
Неизвестные операции & unknown & credit/debit \\
\bottomrule
\end{tabular}
\end{table}

\section{Rule-based классификатор}

\subsection{Принцип работы}

Rule-based (правило-ориентированный) подход является одним из старейших методов классификации текста и сохраняет актуальность в задачах, где:
\begin{itemize}
  \item предметная область хорошо формализована и имеет устойчивую терминологию;
  \item требуется полная интерпретируемость каждого решения;
  \item объём обучающих данных недостаточен для статистических методов \cite{Cohen1995}.
\end{itemize}

Классификатор представляет собой упорядоченный набор правил, применяемых последовательно в порядке убывания специфичности. Более специфичные правила имеют приоритет над более общими. Каждое правило проверяет наличие паттерна в одном или нескольких полях: \texttt{payment\_purpose}, \texttt{counterparty\_name}, \texttt{counterparty\_inn}, а также тип операции \texttt{direction}.

\subsection{Структура правил}

Иерархия правил организована в пять уровней:

\textbf{Уровень 1 — детерминированные идентификаторы.} Операции, однозначно определяемые по ИНН или точному наименованию контрагента. Например, платежи с ИНН налоговой инспекции всегда относятся к категории «налоги и взносы».

\textbf{Уровень 2 — ключевые слова в назначении.} Регулярные выражения по семантически нагруженным токенам:
\begin{itemize}
  \item назначение содержит \texttt{/зарплат|аванс\b|подотчёт/i} → «персонал»;
  \item назначение содержит \texttt{/мука|сахар|меланж|сгущ|молокосодерж|яичн|масло\s+раст/i} → «производственные закупки»;
  \item назначение содержит \texttt{/упаковк|пакет|плёнк|стрейч|лоток|коробк/i} → «упаковка»;
  \item назначение содержит \texttt{/аренд|коммун|электроэнерг|водоснабж|вывоз\s+отходов/i} → «аренда и коммунальные»;
  \item назначение содержит \texttt{/погашение\s+кредит|процент\s+по\s+кредит|аннуитет/i} → «кредиты и займы».
\end{itemize}

\textbf{Уровень 3 — контрагентные словари.} Справочник контрагентов, классифицированных по юридическому лицу: ООО «Меланж» → «производственные закупки», ИП Баскаков С.А. → «аренда и коммунальные», ФНС / ИФНС → «налоги и взносы».

\textbf{Уровень 4 — общие признаки.} Широкие правила для операций, не попавших в предыдущие уровни: кредитовые операции от юридических лиц по умолчанию → «оплаты от клиентов» с низкой уверенностью.

\textbf{Уровень 5 — unknown.} Операции, не подпавшие ни под одно правило.

\subsection{Управление уверенностью}

Каждому правилу присвоен уровень уверенности: высокий (\texttt{high}), средний (\texttt{medium}) или низкий (\texttt{low}). Уверенность высокая, если операция однозначно классифицирована по ИНН или точному паттерну; средняя — при нескольких возможных трактовках; низкая — при классификации по широким правилам. Операции с низкой уверенностью автоматически помечаются флагом \texttt{needs\_review}.

\section{LLM-разметка как дополнительный сигнал}

\subsection{Роль языковых моделей в классификации транзакций}

Большие языковые модели (LLM) демонстрируют высокую эффективность в задачах классификации коротких текстов с нестандартными или многозначными формулировками \cite{brown2020language}. В финансовом контексте LLM способны интерпретировать аббревиатуры, деловой жаргон и неполные описания, которые плохо ловятся регулярными выражениями.

Вместе с тем применение LLM к банковским данным требует особой осторожности с точки зрения конфиденциальности. В данном проекте используется локально развёрнутый LLM-сервер на базе LM Studio, что исключает передачу финансовых данных во внешние сервисы.

\subsection{Архитектура взаимодействия с LLM}

Каждый запрос к модели содержит структурированный prompt следующего вида:

\begin{enumerate}
  \item Системная инструкция с описанием задачи и перечнем допустимых категорий.
  \item Тип операции (\texttt{credit} или \texttt{debit}).
  \item Наименование контрагента и назначение платежа.
  \item Явное требование вернуть ответ строго в формате JSON с полями \texttt{category}, \texttt{subcategory} и \texttt{confidence}.
\end{enumerate}

Ограничение списком допустимых категорий принципиально важно: без него модель склонна к «галлюцинациям» — возврату произвольных категорий, которые невозможно автоматически проверить \cite{ji2023survey}. Формат JSON позволяет автоматически валидировать ответ и отклонять некорректные.

\subsection{Сравнение методов и итоговая схема}

В итоговой схеме оба результата — \texttt{category\_rule} и \texttt{category\_llm} — сохраняются и сравниваются. Логика формирования финальной категории (\texttt{category\_final}):
\begin{itemize}
  \item если rule-based дал результат с уверенностью \texttt{high} — используется rule-based;
  \item если rule-based вернул unknown или \texttt{low} — используется LLM-вариант;
  \item если результаты расходятся при уверенности \texttt{medium} — операция помечается \texttt{needs\_review = True} с пояснением \texttt{review\_reason}.
\end{itemize}

Поле \texttt{review\_reason} содержит текстовое описание причины: «расхождение rule vs LLM», «нет ключевых слов», «новый контрагент», «сумма превышает порог».

\section{Результаты классификации}

\subsection{Покрытие категорий}

По итогам классификации из 4 086 банковских операций:
\begin{itemize}
  \item \textbf{1 операция} отнесена в категорию unknown на сумму 5 793,30 руб., что составляет около \textbf{0,01\%} оборота по сумме;
  \item \textbf{238 операций} (5,8\% по количеству) помечены флагом \texttt{needs\_review}; их суммарная стоимость составляет \textbf{31 379 752 руб.}, что указывает на концентрацию неоднозначных операций в крупных платежах, а не мелких.
\end{itemize}

Практически полное устранение unknown-категории является ключевым техническим результатом: классификация пригодна для построения управленческого cash flow без ручного заполнения пропусков.

\subsection{Ручная валидация}

Для оценки качества подготовлена валидационная выборка из \textbf{500 операций} (файл \texttt{outputs/manual\_validation\_500.xlsx}). Выборка стратифицирована: в неё пропорционально включены операции разных категорий, обоих типов (\texttt{credit}/\texttt{debit}), а также все 238 операций с флагом \texttt{needs\_review}.

В выборке вручную заполняются поля \texttt{category\_manual} и \texttt{subcategory\_manual}. По завершении ручной разметки автоматически рассчитываются метрики согласованности: accuracy, macro F1 и confusion matrix для каждой пары (\texttt{category\_rule} vs \texttt{category\_manual}), (\texttt{category\_llm} vs \texttt{category\_manual}), (\texttt{category\_final} vs \texttt{category\_manual}).

Такой подход к оценке качества согласуется с рекомендациями \cite{pustejovsky2012natural} по использованию inter-annotator agreement для валидации автоматической разметки в отсутствие золотого стандарта.

\subsection{Примеры исправленных паттернов}

В ходе ревью unknown-операций и \texttt{needs\_review}-операций выявлены повторяющиеся паттерны, добавленные в справочник:
\begin{itemize}
  \item ООО «Меланж» — меланж и пастеризованный белок (производственные закупки);
  \item «Тензор», «Бизнес-Софт», «Ростелеком» — цифровые сервисы и связь;
  \item «Эко-Сфера», «Хартия» — вывоз отходов (аренда и коммунальные);
  \item «КЗГ», «Марион», «Сиб Агро» — производственные расходники и оборудование;
  \item «ООО ППР» — договор без предмета; классифицирован как профессиональные услуги с низкой уверенностью и флагом ручной проверки.
\end{itemize}

Итоговый вывод этапа: unknown практически устранён, структура расходов и поступлений прозрачна, данные пригодны для финансового аудита и обучения прогнозных моделей.
